% !TeX root = 01_esercitazione_probabilita_att.tex
\begin{center}
  {\Huge \color{mainblue}\textbf{Esercitazione: Calcolo delle Probabilità}}\\[0.1cm]
  \rule{0.82\linewidth}{0.5pt}
\end{center}

\vspace{0.3cm}

\begin{esercizio}{Il catalogo della biblioteca}
Nella biblioteca di una scuola media ci sono dei libri, catalogati e suddivisi in base al loro genere letterario come mostrato nella tabella seguente:

\begin{center}
\begin{spacing}{1.25}
\begin{tabularx}{0.65\linewidth}{|m{4cm}|>{\centering\arraybackslash}X|}
  \hline
  \textbf{Genere di Libro} & \textbf{Numero di Libri} \\ \hline\hline
  Avventura & 120 \\ \hline
  Gialli & 90 \\ \hline
  Fantasy & 110 \\ \hline
  Scienza & 80 \\ \hline
  Fumetti & 100 \\ \hline
\end{tabularx}
\end{spacing}
\end{center}

Immagina di scegliere un libro a caso dalla biblioteca:

\begin{enumerate}[label=\alph*)]
  \item Verifica tramite un calcolo che nella biblioteca ci siano esattamente $500$ libri in totale.
  \item Qual è la probabilità che il libro estratto a caso sia di genere \textbf{scienza}?
  \item Qual è la probabilità che il libro estratto a caso \textbf{non sia} un fumetto?
  \item Supponiamo che Giacomo decida di scegliere un libro tra quelli di genere \textbf{fantasy}. Sapendo che all'interno del genere fantasy vi è una serie speciale composta da $25$ libri sui draghi, qual è la probabilità che il libro scelto da Giacomo appartenga a questa serie?
\end{enumerate}

\responsespace{3.0cm}

\begin{soluzione}
  \textbf{Passaggi:}
  \begin{enumerate}[label=\alph*)]
    \item Sommiamo i libri di ciascun genere per verificare il totale:
    \[
    \text{Totale} = 120 + 90 + 110 + 80 + 100 = 500 \text{ libri}.
    \]
    La somma è esattamente $500$, che rappresenta il nostro numero di **casi possibili**.
    \item I casi favorevoli per il genere ``scienza'' sono $80$. La probabilità è:
    \[
    P(\text{scienza}) = \frac{\text{casi favorevoli}}{\text{casi possibili}} = \frac{80}{500} = \frac{8}{50} = \frac{4}{25} = 0{,}16 = 16\%
    \]
    \item Per l'evento ``non sia un fumetto'', i casi favorevoli sono tutti i libri tranne i fumetti (ossia $500 - 100 = 400$ libri):
    \[
    P(\text{non fumetto}) = \frac{400}{500} = \frac{4}{5} = 0{,}80 = 80\%
    \]
    \item Poiché sappiamo che Giacomo sceglie *solamente* tra i libri di genere \textbf{fantasy}, lo spazio dei casi possibili si restringe ai soli $110$ libri fantasy.
    I casi favorevoli (i libri sui draghi all'interno della sezione fantasy) sono $25$. La probabilità è:
    \[
    P(\text{draghi} \mid \text{fantasy}) = \frac{25}{110} = \frac{5}{22} \approx 0{,}227 = 22{,}7\%
    \]
  \end{enumerate}
\end{soluzione}
\end{esercizio}

\condnewpage

\begin{esercizio}{La vetrina dei gelati}
Una gelateria prepara coni gelato di diversi gusti per l'inizio del weekend. I gelati pronti in cella frigorifera sono così suddivisi in base al gusto:

\begin{center}
\begin{spacing}{1.25}
\begin{tabularx}{0.65\linewidth}{|m{4cm}|>{\centering\arraybackslash}X|}
  \hline
  \textbf{Gusto} & \textbf{Numero di Gelati} \\ \hline\hline
  Cioccolato & 70 \\ \hline
  Fragola & 40 \\ \hline
  Limone & 30 \\ \hline
  Pistacchio & 50 \\ \hline
  Vaniglia & 50 \\ \hline
\end{tabularx}
\end{spacing}
\end{center}

Un cliente entra e sceglie un cono gelato completamente a caso.

\begin{enumerate}[label=\alph*)]
  \item Verifica tramite un calcolo che la gelateria abbia preparato esattamente $240$ gelati.
  \item Qual è la probabilità che il cliente scelga un gelato al \textbf{pistacchio}?
  \item Qual è la probabilità che the gelato scelto \textbf{non sia} al limone?
  \item Sapendo che il gelato estratto a caso è al \textbf{cioccolato}, qual è la probabilità che sia uno dei $20$ coni ricoperti speciali di granella di nocciole?
\end{enumerate}

\responsespace{3.0cm}

\begin{soluzione}
  \textbf{Passaggi:}
  \begin{enumerate}[label=\alph*)]
    \item Sommiamo i gelati preparati:
    \[
    \text{Totale} = 70 + 40 + 30 + 50 + 50 = 240 \text{ gelati}.
    \]
    La somma è esattamente $240$ (casi possibili).
    \item I casi favorevoli per il gusto ``pistacchio'' sono $50$. La probabilità è:
    \[
    P(\text{pistacchio}) = \frac{50}{240} = \frac{5}{24} \approx 0{,}208 = 20{,}8\%
    \]
    \item Per l'evento ``non sia al limone'', i casi favorevoli sono tutti i coni tranne i $30$ al limone (cioè $240 - 30 = 210$ gelati):
    \[
    P(\text{non limone}) = \frac{210}{240} = \frac{21}{24} = \frac{7}{8} = 0{,}875 = 87{,}5\%
    \]
    \item Poiché sappiamo che il cono scelto è al cioccolato, lo spazio dei casi possibili si restringe ai soli $70$ gelati al cioccolato.
    I casi favorevoli (coni con granella all'interno del cioccolato) sono $20$. La probabilità è:
    \[
    P(\text{granella} \mid \text{cioccolato}) = \frac{20}{70} = \frac{2}{7} \approx 0{,}286 = 28{,}6\%
    \]
  \end{enumerate}
\end{soluzione}
\end{esercizio}

\condnewpage

\begin{esercizio}{Lo scaffale delle magliette}
In un negozio di articoli sportivi ci sono delle magliette di cotone, suddivise per colore come riportato nella seguente tabella:

\begin{center}
\begin{spacing}{1.25}
\begin{tabularx}{0.65\linewidth}{|m{4cm}|>{\centering\arraybackslash}X|}
  \hline
  \textbf{Colore} & \textbf{Numero di Magliette} \\ \hline\hline
  Rosse & 80 \\ \hline
  Blu & 100 \\ \hline
  Verdi & 60 \\ \hline
  Nere & 70 \\ \hline
  Bianche & 50 \\ \hline
\end{tabularx}
\end{spacing}
\end{center}

Un commesso preleva a caso una maglietta dallo scaffale.

\begin{enumerate}[label=\alph*)]
  \item Verifica tramite un calcolo che nello scaffale ci siano esattamente $360$ magliette.
  \item Qual è la probabilità che la maglietta scelta sia \textbf{verde}?
  \item Qual è la probabilità che la maglietta scelta \textbf{non sia} blu?
  \item Sapendo che la maglietta scelta è di colore \textbf{nero}, qual è la probabilità che sia una delle $15$ magliette speciali che presentano il logo dorato stampato sul petto?
\end{enumerate}

\responsespace{3.0cm}

\begin{soluzione}
  \textbf{Passaggi:}
  \begin{enumerate}[label=\alph*)]
    \item Sommiamo le magliette presenti sullo scaffale:
    \[
    \text{Totale} = 80 + 100 + 60 + 70 + 50 = 360 \text{ magliette}.
    \]
    La somma è esattamente $360$ (casi possibili).
    \item I casi favorevoli per il colore ``verde'' sono $60$. La probabilità è:
    \[
    P(\text{verde}) = \frac{60}{360} = \frac{1}{6} \approx 0{,}167 = 16{,}7\%
    \]
    \item Per l'evento ``non sia blu'', i casi favorevoli sono tutte le magliette tranne le $100$ blu (cioè $360 - 100 = 260$ magliette):
    \[
    P(\text{non blu}) = \frac{260}{360} = \frac{26}{36} = \frac{13}{18} \approx 0{,}722 = 72{,}2\%
    \]
    \item Poiché sappiamo che la maglietta prelevata è nera, lo spazio dei casi possibili si riduce alle sole $70$ magliette nere.
    I casi favorevoli (magliette nere con logo dorato) sono $15$. La probabilità è:
    \[
    P(\text{logo dorato} \mid \text{nera}) = \frac{15}{70} = \frac{3}{14} \approx 0{,}214 = 21{,}4\%
    \]
  \end{enumerate}
\end{soluzione}
\end{esercizio}

\condnewpage

\begin{esercizio}{Il lancio di un dado e una moneta}
Vengono lanciati contemporaneamente un comune dado a $6$ facce numerate da $1$ a $6$ e una moneta perfettamente bilanciata. 

Risolvi i seguenti quesiti elencando esplicitamente lo spazio dei casi:

\begin{enumerate}[label=\alph*)]
  \item Elenca tutti i possibili esiti combinati del lancio. Quanti sono in totale?
  \item Qual è la probabilità di ottenere un numero pari sul dado e la faccia Testa (T) sulla moneta?
  \item Qual è la probabilità di ottenere un numero strettamente minore di $3$ sul dadi e la faccia Croce (C) sulla moneta?
  \item Qual è la probabilità di ottenere un numero dispari sul dado (indipendentemente dall'esito della moneta)?
  \item Qual è la probabilità di ottenere Testa (T) sulla moneta \textbf{oppure} di ottenere un numero pari a $6$ sul dado?
\end{enumerate}

\responsespace{3.0cm}

\begin{soluzione}
  \textbf{Passaggi:}
  \begin{enumerate}[label=\alph*)]
    \item Ogni lancio di dado ha 6 esiti e ogni lancio di moneta ha 2 esiti. Spazio dei casi possibili formato da $6 \times 2 = 12$ coppie ordinate:
    \[
    \Omega = \{T1, T2, T3, T4, T5, T6, C1, C2, C3, C4, C5, C6\}
    \]
    I casi possibili sono in totale $12$.
    \item Le coppie con numero pari $\{2,4,6\}$ e Testa (T) sono:
    \[
    \text{Favorevoli} = \{T2, T4, T6\} \implies 3 \text{ casi}
    \]
    La probabilità è:
    \[
    P(\text{pari e T}) = \frac{3}{12} = \frac{1}{4} = 0{,}25 = 25\%
    \]
    \item Un numero strettamente minore di $3$ significa che il dado deve mostrare $1$ o $2$. Abbinati a Croce (C) abbiamo:
    \[
    \text{Favorevoli} = \{C1, C2\} \implies 2 \text{ casi}
    \]
    La probabilità è:
    \[
    P(< 3 \text{ e C}) = \frac{2}{12} = \frac{1}{6} \approx 0{,}167 = 16{,}7\%
    \]
    \item Le facce dispari sul dado sono $\{1,3,5\}$. I casi favorevoli sono tutte le coppie con dadi dispari:
    \[
    \text{Favorevoli} = \{T1, T3, T5, C1, C3, C5\} \implies 6 \text{ casi}
    \]
    La probabilità è:
    \[
    P(\text{dispari}) = \frac{6}{12} = \frac{1}{2} = 0{,}5 = 50\%
    \]
    \item L'evento ``T o 6'' è un'unione di eventi non incompatibili. Contiamo i casi favorevoli:
    - Tutte le coppie con la moneta su Testa: $\{T1, T2, T3, T4, T5, T6\}$ (6 casi).
    - Tutte le coppie con il dado su 6: $\{T6, C6\}$ (il caso $T6$ è già stato contato).
    L'insieme delle coppie favorevoli è $\{T1, T2, T3, T4, T5, T6, C6\}$ (7 casi favorevoli).
    La probabilità è:
    \[
    P(T \text{ oppure } 6) = \frac{7}{12} \approx 0{,}583 = 58{,}3\%
    \]
  \end{enumerate}
\end{soluzione}
\end{esercizio}

\condnewpage

\begin{esercizio}{Il lancio di due dadi}
Si lanciano simultaneamente due dadi regolari a $6$ facce distinguibili: uno di colore \textbf{rosso} e uno di colore \textbf{blu}.

\begin{enumerate}[label=\alph*)]
  \item Qual è la probabilità di ottenere due numeri uguali sulle facce dei due dadi (un ``doppio'')?
  \item Qual è la probabilità che la \textbf{somma} dei due numeri ottenuti sia esattamente pari a $8$?
  \item Qual è la probabilità che compaia \textbf{almeno una volta} la faccia $5$ (sul dado rosso, sul dado blu o su entrambi)?
  \item Qual è la probabilità che il numero mostrato sul dado \textbf{rosso} sia strettamente maggiore di quello mostrato sul dado \textbf{blu}?
\end{enumerate}

\responsespace{3.0cm}

\begin{soluzione}
  \textbf{Passaggi:}
  I casi possibili sono in totale $6 \times 6 = 36$ coppie ordinate del tipo $(R, B)$.
  \begin{enumerate}[label=\alph*)]
    \item I casi favorevoli sono le coppie con facce uguali:
    \[
    \text{Favorevoli} = \{(1,1), (2,2), (3,3), (4,4), (5,5), (6,6)\} \implies 6 \text{ casi}
    \]
    La probabilità è:
    \[
    P(\text{numeri uguali}) = \frac{6}{36} = \frac{1}{6} \approx 0{,}167 = 16{,}7\%
    \]
    \item Identifichiamo le coppie la cui somma delle componenti è pari a 8:
    \[
    \text{Favorevoli} = \{(2,6), (3,5), (4,4), (5,3), (6,2)\} \implies 5 \text{ casi}
    \]
    La probabilità è:
    \[
    P(\text{somma } 8) = \frac{5}{36} \approx 0{,}139 = 13{,}9\%
    \]
    \item L'evento ``almeno un 5'' comprende le coppie in cui la prima componente è 5 (riga del 5: 6 casi) o la seconda è 5 (colonna del 5: 6 casi). Il caso $(5,5)$ è in comune.
    I casi favorevoli sono in totale $6 + 6 - 1 = 11$ casi:
    \[
    \text{Favorevoli} = \{(5,1), (5,2), (5,3), (5,4), (5,5), (5,6), (1,5), (2,5), (3,5), (4,5), (6,5)\}
    \]
    La probabilità è:
    \[
    P(\text{almeno un } 5) = \frac{11}{36} \approx 0{,}306 = 30{,}6\%
    \]
    \item Su 36 casi totali:
    - 6 casi sono di parità ($R = B$, vedi punto a).
    - Nei restanti $30$ casi, per simmetria, la metà delle volte il dado rosso sarà maggiore del blu ($R > B$) e l'altra metà sarà minore ($R < B$).
    Quindi, i casi favorevoli a $R > B$ sono $30 / 2 = 15$ casi (ad esempio $(2,1), (3,1), (3,2)...$):
    \[
    P(R > B) = \frac{15}{36} = \frac{5}{12} \approx 0{,}417 = 41{,}7\%
    \]
  \end{enumerate}
\end{soluzione}
\end{esercizio}

\condnewpage

\begin{esercizio}{Due lanci di una moneta}
Una moneta non truccata viene lanciata in aria due volte consecutivamente.

\begin{enumerate}[label=\alph*)]
  \item Qual è la probabilità di ottenere due volte la faccia Testa (\textbf{due teste})?
  \item Qual è la probabilità di ottenere \textbf{almeno una} croce?
  \item Qual è la probabilità di ottenere \textbf{esattamente} una testa?
\end{enumerate}

\responsespace{3.0cm}

\begin{soluzione}
  \begin{center}
  \begin{tikzpicture}[level distance=1.5cm,
    sibling distance=2.2cm,
    every node/.style={circle, draw=mainblue, fill=softblue, inner sep=2.5pt, thick, font=\bfseries\color{mainblue}},
    edge from parent/.style={draw, gray!80, line width=1pt, ->}]
    
    \node[rectangle, rounded corners=3pt, font=\small] (R) {Inizio}
      child {node {T}
        child {node {T} edge from parent}
        child {node {C} edge from parent}
        edge from parent
      }
      child {node {C}
        child {node {T} edge from parent}
        child {node {C} edge from parent}
        edge from parent
      };
  \end{tikzpicture}
  \end{center}

  \textbf{Passaggi:}
  I possibili esiti dell'esperimento sono formato da $2^2 = 4$ coppie ordinate:
  \[
  \Omega = \{TT, TC, CT, CC\}
  \]
  \begin{enumerate}[label=\alph*)]
    \item L'evento ``due teste'' corrisponde all'esito $\{TT\}$ (1 caso favorevole):
    \[
    P(TT) = \frac{1}{4} = 0{,}25 = 25\%
    \]
    \item L'evento ``almeno una croce'' esclude solo l'esito $TT$. I casi favorevoli sono $\{TC, CT, CC\}$ (3 casi):
    \[
    P(\text{almeno una C}) = \frac{3}{4} = 0{,}75 = 75\%
    \]
    \item L'evento ``esattamente una testa'' comprende le sequenze con una sola T e una C: $\{TC, CT\}$ (2 casi favorevoli):
    \[
    P(\text{esattamente una T}) = \frac{2}{4} = \frac{1}{2} = 0{,}5 = 50\%
    \]
  \end{enumerate}
\end{soluzione}
\end{esercizio}

\condnewpage

\begin{esercizio}{Tre lanci di una moneta}
Una moneta non truccata viene lanciata in aria tre volte consecutivamente.

\begin{enumerate}[label=\alph*)]
  \item Qual è la probabilità di ottenere tre volte Testa (\textbf{tre teste})?
  \item Qual è la probabilità di ottenere \textbf{esattamente} due teste nei tre lanci?
  \item Qual è la probabilità di ottenere \textbf{almeno} una testa nei tre lanci?
  \item Qual è la probabilità che il \textbf{primo lancio} sia croce (indipendentemente dagli altri due)?
\end{enumerate}

\responsespace{3.0cm}

\begin{soluzione}
  \begin{center}
  \begin{tikzpicture}[level distance=1.5cm,
    sibling distance=1.6cm,
    every node/.style={circle, draw=mainblue, fill=softblue, inner sep=2pt, thick, font=\bfseries\color{mainblue}},
    edge from parent/.style={draw, gray!80, line width=1pt, ->}]
    
    \node[rectangle, rounded corners=3pt, font=\small] (R) {Inizio}
      child {node {T}
        child {node {T}
          child {node {T} edge from parent}
          child {node {C} edge from parent}
          edge from parent
        }
        child {node {C}
          child {node {T} edge from parent}
          child {node {C} edge from parent}
          edge from parent
        }
        edge from parent
      }
      child {node {C}
        child {node {T}
          child {node {T} edge from parent}
          child {node {C} edge from parent}
          edge from parent
        }
        child {node {C}
          child {node {T} edge from parent}
          child {node {C} edge from parent}
          edge from parent
        }
        edge from parent
      };
  \end{tikzpicture}
  \end{center}

  \textbf{Passaggi:}
  I possibili esiti dell'esperimento (foglie dell'albero a 3 livelli) sono $2^3 = 8$ terne ordinate:
  \[
  \Omega = \{TTT, TTC, TCT, TCC, CTT, CTC, CCT, CCC\}
  \]
  \begin{enumerate}[label=\alph*)]
    \item L'esito favorevole a ``tre teste'' è solo $\{TTT\}$ (1 caso):
    \[
    P(TTT) = \frac{1}{8} = 0{,}125 = 12{,}5\%
    \]
    \item Le terne che contengono esattamente due T e una C sono:
    \[
    \text{Favorevoli} = \{TTC, TCT, CTT\} \implies 3 \text{ casi}
    \]
    La probabilità è:
    \[
    P(\text{esattamente due T}) = \frac{3}{8} = 0{,}375 = 37{,}5\%
    \]
    \item L'evento ``almeno una testa'' esclude solo l'esito $\{CCC\}$. I casi favorevoli sono tutti gli altri 7 esiti:
    \[
    P(\text{almeno una T}) = \frac{7}{8} = 0{,}875 = 87{,}5\%
    \]
    \item I casi in cui il primo lancio è croce (la terna comincia con la lettera $C$) sono:
    \[
    \text{Favorevoli} = \{CTT, CTC, CCT, CCC\} \implies 4 \text{ casi}
    \]
    \item La probabilità è:
    \[
    P(\text{1° lancio C}) = \frac{4}{8} = \frac{1}{2} = 0{,}5 = 50\%
    \]
  \end{enumerate}
\end{soluzione}
\end{esercizio}
